\sectioncentered*{Заключение}
\addcontentsline{toc}{section}{Заключение}

В курсовой работе разработана гибридная интеллектуальная система игрового ассистента
для виртуальной среды Minecraft, объединяющая онтологическую базу знаний, решатель задач
и большую языковую модель в едином процессе обработки запроса пользователя. Задачи,
поставленные во введении, решены.

Проанализированы предметная область, класс решаемых задач и существующие гибридные
интеллектуальные системы. Установлено, что ассистент, построенный только на языковой
модели, не располагает состоянием среды и не может гарантировать соблюдение её правил, а
символьная схема не понимает свободных формулировок. На этой основе сформированы
функциональные и нефункциональные требования, главное из которых~--~нулевая доля
недопустимых действий независимо от формулировки запроса.

Определены состав и назначение компонентов системы. База знаний представлена
онтологией в формате RDF/OWL, которая собирается автоматически из файлов данных среды и
содержит 25\,287 триплетов, 1522 сущности и 1144 рецепта; её дополняют рукописные слои
запретов, разговорных названий, расположения источников добычи, классов угроз и перечня
объектов древа развития. Решатель задач выполняет обратный вывод по графу способов
получения объектов с учётом инвентаря, проверяет допустимость действий с правом вето,
доказывает выполнимость заданий и этапов кампании, выводит меры защиты от угроз, ступень
и древо развития. Интерфейс включает игровой чат и внутриигровой экран клиента среды, а
также веб-страницы, повторяющие протокол клиента.

Спроектирована архитектура с последовательной схемой интеграции подходов. Языковая
модель находится на входе конвейера, где преобразует свободную реплику в намерение, и на
выходе, где формулирует ответ по выходу решателя; все решения, требующие знаний о
предметной области и гарантий корректности, принимаются между этими точками. Модель
называет объект словами пользователя и не видит исполняемую команду ни на одном шаге,
поэтому недопустимое действие не может быть получено формулировкой запроса. Вопросы за
пределами базы знаний, в том числе о других играх, обрабатываются поиском в сети: ответ
строится по найденным фрагментам с указанием источников, а найденный текст не участвует в
выборе действия.

Реализованы база знаний, резолвер, решатель задач, оркестратор, модуль поиска в сети,
серверная часть интерфейса и клиентская модификация среды, а также оценочный контур для
сравнения конфигураций системы.

Проведено тестирование на компонентном, интеграционном и сквозном уровнях: 603
автоматических теста и трёхкратный прогон размеченного набора из 61 реплики на настоящей
языковой модели. Оценено качество системы в сравнении трёх конфигураций на одном наборе и
одной модели. Гибридная система совпала с разметкой в 56 случаях из 61 и не выполнила ни
одного недопустимого действия, тогда как языковая модель без базы знаний дала 38
совпадений и 16 недопустимых действий, а символьная схема~--~29 и 7. При замене модели на
локальную с четырьмя миллиардами параметров гибрид сохранил 52 совпадения и нулевую долю
недопустимых действий. Время построения плана не превышает 154~мс, 95\,\% запросов,
решаемых по базе знаний, получают ответ быстрее 2~с.

Ожидаемый эффект применения системы состоит в том, что пользователь получает ассистента,
который понимает свободные формулировки, но не выполняет действий, противоречащих
правилам среды, объясняет каждое решение цепочкой применённых правил и строит планы от
фактического состояния пользователя. Предложенная схема применима всюду, где языковая
модель управляет действиями в системе с собственным состоянием и правилами: знания и
проверки вынесены за пределы модели, и модель заменяется изменением конфигурации.

Направлениями дальнейшего развития являются распознавание изображения экрана для
определения игры и объекта, о котором спрашивает пользователь, подключение второй
виртуальной среды наполнением предметной части онтологии без изменения решателя,
пополнение базы знаний способами получения, заданными программным кодом среды,
расширение размеченного набора сценариями с состоянием и голосовое взаимодействие
внутри клиента среды.

Работа выполнена автором индивидуально. Автором проведён анализ предметной области и
существующих систем, спроектированы онтология и рукописные слои знаний, разработаны
алгоритмы решателя и архитектура интеграции компонентов, выполнена программная
реализация системы, построен оценочный контур, проведены тестирование и оценка качества.
